\begin{frame}[fragile]{Immediate Undefined Behaviour}
  Assuming \texttt{c} is \texttt{true} most of the time, we might be tempted to simplify the function by removing the branching as follows:
    \begin{columns}[c]
        \begin{column}{0.44\textwidth}
            \centering
            \begin{lstlisting}[language=C++, frame=single, xleftmargin=4pt, xrightmargin=4pt]
int fun(bool c, int v) {
    int div = 0;
    if (c) {
        div = 42 / v;
    }
    return div;
}
            \end{lstlisting}
        \end{column}

        \begin{column}{0.08\textwidth}
            \centering
            \only<1>{\transform}
            \only<2>{\illegaltransform}
        \end{column}

        \begin{column}{0.44\textwidth}
            \centering
            \begin{lstlisting}[language=C++, frame=single, xleftmargin=4pt, xrightmargin=4pt]
int fun(bool c, int v) {
    int div = 42 / v;
    return c ? div : 0;
}
            \end{lstlisting}
        \end{column}
    \end{columns}
    \uncover<2->{
      However this is \textbf{illegal} since \texttt{fun(false, 0)} will trigger UB after optimization
    }
    \vfill
\pdfpcnote{
**First Reveal**<br/>
Present the transformation as a plausible optimization: c is true most of the
time, so hoisting the division looks reasonable. Stop here before revealing the
problem — let the audience think about whether it is correct.<br/>
**Second Reveal**<br/>
The transformation is illegal: fun(false, 0) would divide by zero after the
rewrite, introducing new immediate UB that did not exist in the original program.
Key rule: a compiler transformation must not introduce UB that was absent in the
source.
}

\end{frame}
